\chapter{Methodology}
\label{ch:methodology}

The study analyses what the providers of health tracking apps tell their users
about privacy harms and protection. This chapter describes how the apps were
chosen. The later steps of the study, from data collection to coding, build on
this sample.

\section{Sampling Frame}
\label{sec:frame}
The study covers Android apps offered in Germany. Android is the larger mobile
platform there, with 72.3\,\% of German mobile page views in September
2026~\parencite{statcounter2026os}. The target is 50 apps. This lies between
the 20 and 64 apps of earlier studies that inspected apps in
detail~\parencite[p.~4]{malki2024exploring}\parencite[p.~5]{hutton2018assessing},
and it is few enough to code every policy and to inspect every app by hand.
The apps come from two groups: 49 Google Play apps selected by popularity and
the preinstalled health tracker of the leading Android phone vendor
(Section~\ref{sec:ranking}).

\section{Search}
\label{sec:search}
\paragraph{Keywords.} The health tracking domain follows the health categories
of Hutton et al.\ (cycling, diet, exercise, heart rate, mood, running, sleep,
step count and weight)~\parencite{hutton2018assessing} and the reproductive
health apps studied by Malki et al.\ (menstrual cycle, fertility and
pregnancy)~\parencite{malki2024exploring}. The 15 terms of Hutton et al.\ and
the 12 terms of Malki et al.\ are searched in English and in one German
translation each. A German word is searched only once when it is the same as
the English term or translates two English terms. This gives 49 queries. The
check with well-known apps below may add further queries; the final list is in
Table~\ref{tab:keywords}.

\paragraph{Search.} Google Play is searched in its German storefront, logged
out, with German as the interface language. For each query, the script keeps
up to the first 30 apps of the result page, in the order shown; some queries
return fewer. It then reads the store page of every app found, which gives the
developer, the category, the price, the description and the number of written
reviews. The script is listed in Appendix~\ref{app:script}.

\paragraph{Check with well-known apps.} A sample of popular health trackers
should contain the apps that most people know. Before the search, a list of
25 well-known health trackers was fixed, covering running and cycling,
nutrition, sleep, mood, the menstrual cycle and the companion apps of
wearables and scales; apps that Google Play does not offer in Germany were
left out. Each listed app should appear among the results of at least one
query. For a missing app, a generic term for its function or device was added
in English and German and the search was repeated. Only terms that found a
missing app were kept, with their translation. Brand names were never used,
because they would favour single providers. The check stops when no further
generic term finds a missing app. It tests whether the queries reach the
market; it does not force any app into the sample.

\section{Eligibility}
\label{sec:eligibility}
An app is eligible if it meets all inclusion rules and none of the exclusion
rules in Table~\ref{tab:eligibility}. Rules I1.1 to I1.4 together define a
health tracker, and all four must hold. The rules are checked on the German
Google Play page of the app: its title, developer, category, price and
description.

\begin{table}[H]
\centering
\begin{singlespace}
\footnotesize
\caption[Eligibility rules]{Eligibility rules, checked on the German Google Play page.}
\label{tab:eligibility}
\begin{tabularx}{\textwidth}{@{}l>{\raggedright\arraybackslash}X@{\hspace{12pt}}l>{\raggedright\arraybackslash}X@{}}
\toprule
\multicolumn{2}{@{}l}{\textbf{Include if \dots}} & \multicolumn{2}{l@{}}{\textbf{Exclude if the app \dots}} \\
\midrule
I1.1 & the user enters data, or a sensor or paired device records them & X1 & shows exercises or articles but records nothing \\
I1.2 & the data describe the user's own health or body & X2 & is a game, a shop or a general-purpose app \\
I1.3 & entries are kept and can be viewed again over time & X3 & tracks a child or another person instead of the user \\
I1.4 & the data fall into one of the searched categories & X4 & records only other data, such as spending or time \\
I5 & the app is offered to the general public & X5 & asks for input once and keeps no result \\
I6 & the app is free to install (in-app purchases are allowed) & X6 & has no tracking or predictive feature \\
I7 & the app is available in German or English & X7 & is limited to members, patients or employees of an organisation \\
 & & X8 & is a DiGA, a digital health application reimbursed in Germany \\
\bottomrule
\end{tabularx}
\end{singlespace}
\end{table}

\section{Ranking and Selection}
\label{sec:ranking}
\paragraph{Popularity.} The apps are ranked by their popularity in Germany,
measured as the number of written reviews that Google Play returns for the
German storefront. Download counts are only shown as worldwide ranges, so they
cannot rank apps in the German market. Malki et al.\ also ranked Google Play
apps by their number of reviews~\parencite[p.~4]{malki2024exploring}. The
review count is a proxy for popularity, not a count of users.

\paragraph{Selection.} All apps found by the search are sorted by their review
count, from the highest to the lowest; ties are broken by the package name.
Apps whose store page cannot be read are left out. The apps are then screened
in this order, and every eligible app enters the sample until 49 apps are
selected. Apps from the same developer count as separate apps, because each
has its own store page and interface.

\paragraph{Preinstalled app.} Some health apps come with the phone. They reach users without a download, so
their review counts understate how many people have them. The sample therefore
adds the preinstalled health tracker of the vendor with the largest share of
German mobile page views from Android
devices~\parencite{statcounter2026vendor}. The app must be made by the vendor
itself, be installed on the vendor's current phones without a download and
meet the eligibility rules. It is screened on its own and skipped in the
ranking, so that it is not counted twice.
